
\part{Application thorn writing}

%%%%%%%%%%%%%%%%%%%%%%%%%%%%%%%%%%%%%%%%%%%%%%%%%%%%%%%%%%%%%%%%%%%%%%%
%%%%%%%%%%%%%%%%%%%%%%%%%%%%%%%%%%%%%%%%%%%%%%%%%%%%%%%%%%%%%%%%%%%%%%%

\chapter{Thorn concepts} 

      (This section has to contain enough explanation to make the rest of
       the writers guide readable the first time through)
(The hardest bugs to find are
those arising from plausible but incorrect assumptions about the
behavior of someone else's thorn.)
      a) Again probably emphasize collaboration, what are thorns,
         packages, how to share them.
      b) Things to think about before you start programming:
             Language, read all the documentation, emphasize use of
             standard supported Cactus infrastructure
      c) Available data types
         i)   Scalars
         ii)  Arrays and GFs
         iii) Groups
      d) Ghost zones and parallelism
      e) Understanding the RFR concept
      f) Understanding the GH concept

%%%%%%%%%%%%%%%%%%%%%%%%%%%%%%%%%%%%%%%%%%%%%%%%%%%%%%%%%%%%%%%%%%%%%%%
%%%%%%%%%%%%%%%%%%%%%%%%%%%%%%%%%%%%%%%%%%%%%%%%%%%%%%%%%%%%%%%%%%%%%%%

\chapter{Anatomy of a thorn}

Each thorn must be in a package. Packages are fairly informal. If you
want, you can just make your own package directory, and stick your new
thorn into it.

      a) Creating a thorn (gmake new thorn)
      b) Thorn directory structure
      c) Thorn make files
      d) Thorn configuration
         (i)   Parameters (param.ccl)
         (ii)  Variables (interface.ccl)
         (iii) Scheduler (schedule.ccl)

%%%%%%%%%%%%%%%%%%%%%%%%%%%%%%%%%%%%%%%%%%%%%%%%%%%%%%%%%%%%%%%%%%%%%%%
%%%%%%%%%%%%%%%%%%%%%%%%%%%%%%%%%%%%%%%%%%%%%%%%%%%%%%%%%%%%%%%%%%%%%%%

\chapter{Putting code into your thorn}

\section{What the F***h provides}

Header files, macros ...

\subsection{IO}

To allow flexible IO, the flesh itself does not provide 
any output routines, however it provides a mechanism for 
thorns to register different routines as IO methods. For 
details of writing IO thorns see ????. Application thorns
can interact with the different IO methods through the following
function calls:

{\t
\begin{verbatim}

#include ``CactusIOFunctions.h''
int CCTK_OutputGH(cGH *GH);

#include ``cctk.h''
CCTK_OutputGH(GH)
    CCTK_POINTER GH
\end{verbatim}
}

\vskip .25cm

This call loops over all registered IO methods, calling 
the routine that each method has registered for {\t OutputGH}.
The expected behaviour of any methods {\t OutputGH} is to
loop over all GH variables outputting them if the method 
contains appropriate routines (that is, not all methods will 
supply routines to output all different types of variables) 
and if the method decides it is an appropriate time to 
output. 


{\t
\begin{verbatim}

#include ``CactusIOFunctions.h''
int CCTK_OutputVarAsByMethod(cGH *GH, const char *varname, const char *alias, const char *methodname);

#include ``cctk.h''
CCTK_OutputVarAsByMethod(GH,varname,alias,methodname)
    CCTK_POINTER GH
    char* varname
    char* alias
    char* methodname
\end{verbatim}
}

Output a variable {\t varname} using the method {\t methodname} if it is 
registered. Uses {\t alias} as the name of the variable for the purpose
of constructing a filename. The output should take place if at all possible,
if the appropriate file exists the data is appended, otheriwise a new
file is created.


{\t
\begin{verbatim}

#include ``CactusIOFunctions.h''
int CCTK_OutputVarByMethod(cGH *GH, const char *varname, const char *methodname);

#include ``cctk.h''
CCTK_OutputVarByMethod(GH,varname,methodname)
    CCTK_POINTER GH
    char* varname
    char* methodname
\end{verbatim}
}

Output a variable {\t varname} using the method {\t methodname} if it is 
registered. The output should take place if at all possible,
if the appropriate file exists the data is appended, otherwise a new
file is created.

{\t
\begin{verbatim}

#include ``CactusIOFunctions.h''
int CCTK_OutputVarAs(cGH *GH, const char *varname, const char *alias);

#include ``cctk.h''
CCTK_OutputVarAs(GH,varname,alias)
    CCTK_POINTER GH
    char* varname
    char* alias
\end{verbatim}
}

Output a variable {\t varname} looping over all registered methods. 
The output should take place if at all possible,
if the appropriate file exists the data is appended, otherwise a new
file is created. Uses {\t alias} as the name of the variable for the purpose
of constructing a filename.

{\t
\begin{verbatim}

#include ``CactusIOFunctions.h''
int CCTK_OutputVar(cGH *GH, const char *varname);

#include ``cctk.h''
CCTK_OutputVarAs(GH,varname)
    CCTK_POINTER GH
    char* varname
\end{verbatim}
}

Output a variable {\t varname} looping over all registered methods. 
The output should take place if at all possible,
if the appropriate file exists the data is appended, otherwise a new
file is created.



\section{Argument lists and parameters}

\section{A First Example (Baloney)}

\section{Programming language differences}

\section{A more complex example (WaveToy)}

\section{Error handling, Warnings and Code Termination}
      
There are two CCTK commands to use for stopping the code
from within your thorn:
\begin{itemize}
\item{} To shut the code down cleanly, use
{\t
\begin{verbatim}
CCTK_Stop(pointer GH, INTEGER return_code)
CCTK_Stop(cGH *GH, int return_code);
\end{verbatim}
}
\item{} To shut the code down more violently, use
{\t
\begin{verbatim}
CCTK_Abort(pointer GH, INTEGER return_code)
CCTK_Abort(cGH *GH, int return_code);
\end{verbatim}
}
\end{itemize}
In both cases, an error code should be returned indicating the
error that you trapped which precipitated the code shutdown.
The error codes are detailed in the following table:
{\bf ACTION: Fill in the table QUERY}.

\section{Calls between different programming languages}
\subsection{Calling C routines from FORTRAN}
\subsection{Calling FORTRAN routines from C}

\section{Programming Style Guidelines and Recommendations}

\section{Adding a test suite}

\section{Sharing your thorns/packages with others}

%%%%%%%%%%%%%%%%%%%%%%%%%%%%%%%%%%%%%%%%%%%%%%%%%%%%%%%%%%%%%%%%%%%%%%%


